\section*{Chapitre \ref{ch:dv:monte-carlo-pricers-in-practice} --- Pricers de Monte-Carlo en pratique}

\begin{solution}{exo:dv:monte-carlo-pricers-in-practice:1}
Le prix de l'option bermudéenne est la borne supérieure, sur l'ensemble des règles d'arrêt, de l'espérance de la valeur d'exercice actualisée. La politique ajustée est une règle d'arrêt, donc son
espérance est au plus égale au prix. Sur des trajectoires indépendantes de l'ajustement, la politique est fixée avant le tirage des trajectoires, et la moyenne est un
estimateur sans biais de la valeur de cette politique, qui est une borne inférieure.
\end{solution}

\begin{solution}{exo:dv:monte-carlo-pricers-in-practice:2}
$S_T=S_0\exp\bigl((r-\tfrac12\sigma^2)T+\sigma\sqrt T Z\bigr)$, donc $\partial S_T/\partial S_0=S_T/S_0$. Le payoff $(S_T-K)^+$ a pour dérivée
$\mathbf 1\{S_T>K\}$ sauf en $S_T=K$, événement de probabilité nulle. La règle de dérivation en chaîne donne alors $e^{-rT}\mathbf 1\{S_T>K\}S_T/S_0$ sur chaque trajectoire, et sa
moyenne est le delta.
\end{solution}

\begin{solution}{exo:dv:monte-carlo-pricers-in-practice:3}
Le payoff de l'\omterm{def:dv:barriers-and-digitals:digital}{option digitale} est constant de part et d'autre du strike, donc sa dérivée est nulle presque partout sur les trajectoires. Le delta provient entièrement de la masse de probabilité
qui traverse le strike, et que la méthode trajectorielle, qui dérive sous l'espérance, ne voit pas~; l'interversion de la dérivation et de l'espérance échoue
pour un payoff discontinu. Le rapport de vraisemblance dérive la densité de $S_T$, qui est lisse, et conserve le payoff tel quel.
\end{solution}

\begin{solution}{exo:dv:monte-carlo-pricers-in-practice:4}
$x=\ln S_T$ suit une loi normale d'espérance $\mu=\ln S_0+(r-\tfrac12\sigma^2)T$ et d'écart-type $\sigma\sqrt T$. Sa log-densité vaut
$-\ln(\sigma\sqrt T)-\tfrac12Z^2+\text{constante}$ avec $Z=(x-\mu)/(\sigma\sqrt T)$. En dérivant par rapport à $\sigma$, à $x$ fixé~:
$\partial Z/\partial\sigma=-Z/\sigma-(\partial\mu/\partial\sigma)/(\sigma\sqrt T)=-Z/\sigma+\sqrt T$, d'où le score vaut
$-1/\sigma-Z(-Z/\sigma+\sqrt T)=(Z^2-1)/\sigma-Z\sqrt T$.
\end{solution}

\begin{solution}{exo:dv:monte-carlo-pricers-in-practice:5}
Conditionnellement à la trajectoire externe, les estimations internes sont les vraies valeurs de continuation plus un bruit $\varepsilon_j$ d'espérance nulle. La borne calcule la
moyenne de $\max_j(a_j+\varepsilon_j)$, où $a_j$ est la valeur que prendrait la borne avec des valeurs exactes. Le maximum est convexe, donc, par l'inégalité de Jensen,
$\E[\max_j(a_j+\varepsilon_j)]\ge\max_j a_j$. Le bruit ne peut que relever la borne, et davantage de trajectoires internes réduisent l'excès, comme le montrent les
valeurs 6.134, 6.054 et 6.037 du put.
\end{solution}

\begin{solution}{exo:dv:monte-carlo-pricers-in-practice:6}
Les coordonnées pseudo-aléatoires sont indépendantes et identiquement distribuées. Le choix de la normale qui pilote chaque morceau de la trajectoire ne modifie pas
la loi de l'estimateur, donc le pont brownien n'y change rien. Les coordonnées de suite à discrépance faible ne sont pas toutes aussi bonnes~: les premières sont les mieux réparties, et en grande
dimension les suivantes sont mal réparties deux à deux. Le pont brownien concentre l'essentiel de la variance de la trajectoire sur les premières coordonnées, ce qui abaisse la
dimension effective de l'intégrande.
\end{solution}

\begin{solution}{exo:dv:monte-carlo-pricers-in-practice:7}
Avec 2\,000 trajectoires d'apprentissage, moyennées sur 40 ensembles d'apprentissage, l'ajustement et la valorisation sur les mêmes trajectoires donnent 18.74~; les mêmes politiques appliquées à 100\,000
trajectoires indépendantes donnent 18.50. L'écart vaut $+0.25$ (erreur type 0.06), et le prix en échantillon dépasse même la \omterm{def:dv:monte-carlo-pricers-in-practice:dual}{borne supérieure duale} 18.713. La
régression a vu le futur de chaque trajectoire, si bien que sur ces trajectoires elle exerce avec clairvoyance~: le biais est à la hausse. Avec 200\,000 trajectoires d'apprentissage et neuf
régresseurs, le biais passe sous le bruit.
\end{solution}

\begin{solution}{exo:dv:monte-carlo-pricers-in-practice:8}
L'erreur type ne mesure que le bruit de l'estimation de la borne inférieure. Le prix est biaisé à la baisse par la sous-optimalité de la politique, et l'erreur type
ne voit pas ce biais~: avec la petite base, l'erreur type vaut 0.039, mais les bornes sont séparées de 0.668. Deux centimes de précision exigent aussi la borne duale (et une
vérification de la discrétisation en temps lorsque le modèle en comporte une).
\end{solution}

\begin{solution}{pb:dv:monte-carlo-pricers-in-practice:1}
\textbf{1.} 18.649, erreur type 0.039.
\textbf{2.} La valorisation sur les trajectoires d'ajustement est biaisée à la hausse (exercice~7~: $+0.25$ avec 2\,000 trajectoires d'apprentissage)~; des trajectoires indépendantes font du prix une borne inférieure.
\textbf{3.} 18.137.
\textbf{4.} Elle ne voit que le meneur. Lorsque le second prix est proche du premier, l'option sur le maximum vaut davantage si on la conserve, et la petite base ne peut
distinguer ces états~; elle exerce trop tôt là, et trop tard lorsque le meneur est loin devant.
\textbf{5.} 15.57 (erreur type 0.07)~: une borne inférieure pour l'option bermudéenne. La \omterm{def:dv:american-options-and-early-exercise:premium}{prime d'exercice anticipé}, d'environ 3.08, provient du rendement de dividende de 10~\%, que le
détenteur abandonne en attendant.
\textbf{6.} 18.713, erreur type 0.019.
\textbf{7.} 0.064, soit 0.3~\% du prix.
\textbf{8.} 18.805~; l'écart vaut 0.668, dix fois plus large.
\textbf{9.} 6.134, 6.054 et 6.037 face au 6.0326 de l'arbre~: l'excès par rapport à l'arbre est divisé par environ cinq à chaque fois que les trajectoires internes sont multipliées par quatre.
\textbf{10.} Chaque nœud externe requiert sa propre simulation interne~: $1\,000\times1\,000$ trajectoires internes partant de chacune des neuf dates, prolongées jusqu'à la dernière date,
soit 45 millions de pas de trajectoire, contre 1.8 million pour les 200\,000 trajectoires de la borne inférieure, soit 25 fois plus.
\textbf{11.} Par la méthode trajectorielle, en gardant la politique ajustée fixe~: à la politique optimale, la sensibilité de la valeur à la frontière est nulle au premier ordre, si bien que la fixer coûte
peu. Ou bien faire un bump avec des nombres aléatoires communs et la même politique.
\textbf{12.} La différence de deux estimations indépendantes a pour variance la somme des deux variances, divisée par le carré du bump, énorme pour un petit bump. Réajuster la
politique sur de nouvelles trajectoires ajoute un saut qui lui est propre.
\textbf{13.} Toutes les sensibilités (trois spots, trois volatilités, corrélations, taux) de chaque transaction pour le coût de quelques valorisations, indépendamment de leur nombre.
\textbf{14.} La trajectoire est un point de dimension 27 (trois actifs à neuf dates)~; un pont brownien par actif et un ordonnancement par composantes principales entre les actifs abaissent la
dimension effective. La décision d'exercice est une discontinuité, ce qui limite le gain.
\textbf{15.} Aucune en temps~: les prix lognormaux sont simulés exactement aux dates d'exercice. Une \omterm{def:dv:local-volatility:model}{volatilité locale} ou stochastique exigerait de petits pas entre les
dates et introduirait un biais.
\textbf{16.} Le desk est vendeur, et le détenteur peut mieux exercer que la politique du desk, si bien que le passif peut atteindre la borne supérieure. Coter autour du
milieu des bornes et comptabiliser la borne inférieure avec une provision pour l'écart.
\textbf{17.} L'écart, plus environ deux erreurs types des bornes, soit $0.064+0.086\approx0.15$ pour 100 de notionnel~; avec la petite base, ce serait environ 0.8.
\textbf{18.} Lorsque l'écart est dans le bruit statistique des bornes~: 0.064 contre deux erreurs types combinées de 0.086, si bien que la base riche l'a atteint.
\textbf{19.} Base riche~: 18.649 contre 18.713, un écart de 0.064 (0.3~\% du prix). Petite base~: 18.137 contre 18.805, un écart de 0.668, dix fois plus large.
\textbf{20.} L'erreur maximale que peut comporter le prix du fait de la politique d'exercice, à laquelle s'ajoute le biais à la hausse de la simulation interne.
\end{solution}

\begin{solution}{iq:dv:monte-carlo-pricers-in-practice:1}
Simuler les trajectoires vers l'avant. En remontant depuis la dernière date d'exercice, régresser les flux actualisés réalisés sur des fonctions de base de l'état, en utilisant
les trajectoires dans la monnaie, et exercer là où la valeur d'exercice dépasse la continuation ajustée~; propager les flux réalisés vers l'arrière. Puis appliquer la politique à de nouvelles
trajectoires~: ce prix est une borne inférieure.
\iqlookfor{régression dans la monnaie, flux réalisés, trajectoires de valorisation indépendantes, borne inférieure.}
\end{solution}

\begin{solution}{iq:dv:monte-carlo-pricers-in-practice:2}
Par dualité~: le prix est au plus égal à $\E[\max_j(h_j-M_j)]$ pour toute martingale $M$ telle que $M_0=0$, avec égalité pour la bonne martingale. Andersen--Broadie construisent $M$ à partir
d'une politique de borne inférieure avec une simulation imbriquée de ses valeurs de continuation. L'écart entre les bornes mesure la qualité de la politique.
\iqlookfor{formule duale, simulation imbriquée, écart.}
\end{solution}

\begin{solution}{iq:dv:monte-carlo-pricers-in-practice:3}
La méthode trajectorielle dérive le payoff le long de chaque trajectoire~: variance faible, mais elle exige un payoff continu par rapport au paramètre~; elle échoue pour les digitales et les barrières.
Le rapport de vraisemblance dérive la densité~: tout payoff convient, mais la variance est plus élevée (6.5 fois pour le delta d'un call dans le chapitre), et il faut une densité connue,
peu commode après de nombreux pas. Le lissage ou le conditionnement combine les deux méthodes.
\iqlookfor{continuité contre densité, et compromis de variance.}
\end{solution}

\begin{solution}{iq:dv:monte-carlo-pricers-in-practice:4}
Le \omterm{def:dv:monte-carlo-pricers-in-practice:qe}{schéma quadratique-exponentiel} d'Andersen pour la variance, avec la variance intégrée prise aux deux bornes du pas pour le log-prix, et une
correction de martingale si nécessaire. Euler rend la variance négative lorsque la condition de Feller n'est pas satisfaite, et la troncature laisse un biais qui décroît lentement (0.86 avec quatre
pas par an dans le chapitre, contre QE dans le bruit statistique).
\iqlookfor{QE, et pourquoi la troncature d'Euler est biaisée.}
\end{solution}

\begin{solution}{iq:dv:monte-carlo-pricers-in-practice:5}
Fixer les graines et utiliser les mêmes nombres aléatoires pour les calculs de base et ceux avec bump~; garder la politique d'exercice fixe~; remplacer les bumps par des \omterm{def:dv:monte-carlo-pricers-in-practice:pathwise}{grecques trajectorielles}, des \omterm{def:dv:monte-carlo-pricers-in-practice:lr}{grecques par rapport de
vraisemblance} ou des \omterm{def:dv:greeks-and-the-hedging-pnl:greeks}{grecques} adjointes~; lisser les payoffs discontinus~; ajouter des variables de contrôle~; vérifier que c'est le bruit, et non le marché, qui fait bouger les chiffres.
\iqlookfor{nombres aléatoires communs et \omterm{def:dv:greeks-and-the-hedging-pnl:greeks}{grecques} directes.}
\end{solution}

\begin{solution}{iq:dv:monte-carlo-pricers-in-practice:6}
Simuler les trajectoires vers l'avant et stocker les états~; parcourir ensuite chaque trajectoire à rebours, en propageant la dérivée du payoff par rapport à chaque état (la
méthode adjointe) et en accumulant les sensibilités à chaque entrée. Toutes les \omterm{def:dv:greeks-and-the-hedging-pnl:greeks}{grecques} coûtent un petit multiple d'une valorisation, quel que soit leur nombre~; le prix à payer est la mémoire pour
les trajectoires stockées (ou un recalcul par checkpointing), et la nécessité, pour une méthode trajectorielle, de payoffs continus.
\iqlookfor{passe arrière, coût indépendant du nombre d'entrées, mémoire.}
\end{solution}
